% 原创TikZ重绘；层次布局参考Bishop与Bishop（2024）第12章。
\begin{tikzpicture}[x=1mm,y=1mm,font=\figurefont\footnotesize,
 every node/.style={text=NNInk},
 stage/.style={nn operator,minimum height=9mm,inner sep=2pt},
 norm/.style={stage,minimum height=7mm},
 wire/.style={nn flow,draw=NNWire,rounded corners=1.5mm},
 aux/.style={draw=NNLine,line width=.85pt,densely dashed},
 sum/.style={nn pointwise,minimum size=5mm}]

\node[nn heading] (encTitle) at (33,140.25) {编码器块 ×6};
\node[nn heading] (decTitle) at (127,140.25) {解码器块 ×6};
\node[] (oStok0) at (15,-4.25) {猫};
\node[inner sep=0pt,minimum width=12mm,minimum height=3mm] (oSemb0) at (15,5.95) {};
\filldraw[draw=NNInput,fill=NNInput!18,line width=.5pt] (9.0,4.675) rectangle ++(3.0,2.55);
\filldraw[draw=NNInput,fill=NNInput!18,line width=.5pt] (12.0,4.675) rectangle ++(3.0,2.55);
\filldraw[draw=NNInput,fill=NNInput!18,line width=.5pt] (15.0,4.675) rectangle ++(3.0,2.55);
\filldraw[draw=NNInput,fill=NNInput!18,line width=.5pt] (18.0,4.675) rectangle ++(3.0,2.55);
\draw[wire] (oStok0.north)--(oSemb0.south);
\node[] (oStok1) at (33,-4.25) {追};
\node[inner sep=0pt,minimum width=12mm,minimum height=3mm] (oSemb1) at (33,5.95) {};
\filldraw[draw=NNInput,fill=NNInput!18,line width=.5pt] (27.0,4.675) rectangle ++(3.0,2.55);
\filldraw[draw=NNInput,fill=NNInput!18,line width=.5pt] (30.0,4.675) rectangle ++(3.0,2.55);
\filldraw[draw=NNInput,fill=NNInput!18,line width=.5pt] (33.0,4.675) rectangle ++(3.0,2.55);
\filldraw[draw=NNInput,fill=NNInput!18,line width=.5pt] (36.0,4.675) rectangle ++(3.0,2.55);
\draw[wire] (oStok1.north)--(oSemb1.south);
\node[] (oStok2) at (51,-4.25) {狗};
\node[inner sep=0pt,minimum width=12mm,minimum height=3mm] (oSemb2) at (51,5.95) {};
\filldraw[draw=NNInput,fill=NNInput!18,line width=.5pt] (45.0,4.675) rectangle ++(3.0,2.55);
\filldraw[draw=NNInput,fill=NNInput!18,line width=.5pt] (48.0,4.675) rectangle ++(3.0,2.55);
\filldraw[draw=NNInput,fill=NNInput!18,line width=.5pt] (51.0,4.675) rectangle ++(3.0,2.55);
\filldraw[draw=NNInput,fill=NNInput!18,line width=.5pt] (54.0,4.675) rectangle ++(3.0,2.55);
\draw[wire] (oStok2.north)--(oSemb2.south);
\node[stage,draw=NNInput,fill=NNInput!18,minimum width=54mm,minimum height=9mm] (oSembedding) at (33,19.55) {叠加正弦位置编码};
\draw[wire] (oSemb0.north)--(15,15.3);
\draw[wire] (oSemb1.north)--(33,15.3);
\draw[wire] (oSemb2.north)--(51,15.3);
\node[] (oSinput) at (33,30.6) {$\bm H_{\mathrm s}^{(0)}$};
\draw[wire] (oSembedding.north)--(oSinput.south);
\node[] (oTtok0) at (109,-4.25) {BOS};
\node[inner sep=0pt,minimum width=12mm,minimum height=3mm] (oTemb0) at (109,5.95) {};
\filldraw[draw=NNInput,fill=NNInput!18,line width=.5pt] (103.0,4.675) rectangle ++(3.0,2.55);
\filldraw[draw=NNInput,fill=NNInput!18,line width=.5pt] (106.0,4.675) rectangle ++(3.0,2.55);
\filldraw[draw=NNInput,fill=NNInput!18,line width=.5pt] (109.0,4.675) rectangle ++(3.0,2.55);
\filldraw[draw=NNInput,fill=NNInput!18,line width=.5pt] (112.0,4.675) rectangle ++(3.0,2.55);
\draw[wire] (oTtok0.north)--(oTemb0.south);
\node[] (oTtok1) at (127,-4.25) {cat};
\node[inner sep=0pt,minimum width=12mm,minimum height=3mm] (oTemb1) at (127,5.95) {};
\filldraw[draw=NNInput,fill=NNInput!18,line width=.5pt] (121.0,4.675) rectangle ++(3.0,2.55);
\filldraw[draw=NNInput,fill=NNInput!18,line width=.5pt] (124.0,4.675) rectangle ++(3.0,2.55);
\filldraw[draw=NNInput,fill=NNInput!18,line width=.5pt] (127.0,4.675) rectangle ++(3.0,2.55);
\filldraw[draw=NNInput,fill=NNInput!18,line width=.5pt] (130.0,4.675) rectangle ++(3.0,2.55);
\draw[wire] (oTtok1.north)--(oTemb1.south);
\node[] (oTtok2) at (145,-4.25) {chases};
\node[inner sep=0pt,minimum width=12mm,minimum height=3mm] (oTemb2) at (145,5.95) {};
\filldraw[draw=NNInput,fill=NNInput!18,line width=.5pt] (139.0,4.675) rectangle ++(3.0,2.55);
\filldraw[draw=NNInput,fill=NNInput!18,line width=.5pt] (142.0,4.675) rectangle ++(3.0,2.55);
\filldraw[draw=NNInput,fill=NNInput!18,line width=.5pt] (145.0,4.675) rectangle ++(3.0,2.55);
\filldraw[draw=NNInput,fill=NNInput!18,line width=.5pt] (148.0,4.675) rectangle ++(3.0,2.55);
\draw[wire] (oTtok2.north)--(oTemb2.south);
\node[stage,draw=NNInput,fill=NNInput!18,minimum width=54mm,minimum height=9mm] (oTembedding) at (127,19.55) {叠加正弦位置编码};
\draw[wire] (oTemb0.north)--(109,15.3);
\draw[wire] (oTemb1.north)--(127,15.3);
\draw[wire] (oTemb2.north)--(145,15.3);
\node[] (oTinput) at (127,30.6) {$\bm H_{\mathrm t}^{(0)}$};
\draw[wire] (oTembedding.north)--(oTinput.south);
\draw[nn frame] (3,36.55) rectangle (66,121.55);
\draw[nn frame] (96,36.55) rectangle (162,121.55);
\node[stage,minimum width=45mm,minimum height=9mm] (e0) at (33,49.3) {双向多头自注意力};
\draw[wire] (oSinput.north)--(e0.south);
\node[norm,minimum width=45mm,minimum height=7mm] (e1) at (33,64.6) {Add \& Norm};
\draw[wire] (e0.north)--(e1.south);
\node[stage,minimum width=45mm,minimum height=9mm] (e2) at (33,85.85) {逐位置FFN};
\draw[wire] (e1.north)--(e2.south);
\node[norm,minimum width=45mm,minimum height=7mm] (e3) at (33,102.0) {Add \& Norm};
\draw[wire] (e2.north)--(e3.south);
\draw[wire] (33,39.1)--(61,39.1)--(61,64.6)--(e1.east);
\fill[NNWire] (33,39.1) circle (.55);
\draw[wire] (33,73.1)--(61,73.1)--(61,102.0)--(e3.east);
\fill[NNWire] (33,73.1) circle (.55);
\node[] (sourcefinal) at (33,128.35) {$\bm H_{\mathrm s}$};
\draw[wire] (e3.north)--(sourcefinal.south);
\node[stage,minimum width=45mm,minimum height=9mm] (d0) at (127,46.75) {因果多头自注意力};
\draw[wire] (oTinput.north)--(d0.south);
\node[norm,minimum width=45mm,minimum height=7mm] (d1) at (127,58.65) {Add \& Norm};
\draw[wire] (d0.north)--(d1.south);
\node[stage,minimum width=45mm,minimum height=9mm] (d2) at (127,75.65) {交叉注意力};
\draw[wire] (d1.north)--(d2.south);
\node[norm,minimum width=45mm,minimum height=7mm] (d3) at (127,87.55) {Add \& Norm};
\draw[wire] (d2.north)--(d3.south);
\node[stage,minimum width=45mm,minimum height=9mm] (d4) at (127,102.85) {逐位置FFN};
\draw[wire] (d3.north)--(d4.south);
\node[norm,minimum width=45mm,minimum height=7mm] (d5) at (127,114.75) {Add \& Norm};
\draw[wire] (d4.north)--(d5.south);
\draw[wire] (127,39.1)--(155,39.1)--(155,58.65)--(d1.east);
\fill[NNWire] (127,39.1) circle (.55);
\draw[wire] (127,65.45)--(155,65.45)--(155,87.55)--(d3.east);
\fill[NNWire] (127,65.45) circle (.55);
\draw[wire] (127,93.5)--(155,93.5)--(155,114.75)--(d5.east);
\fill[NNWire] (127,93.5) circle (.55);
\node[stage,draw=NNOutput,fill=NNOutput!18,minimum width=45mm,minimum height=8mm] (lmout) at (127,128.35) {词表投影＋Softmax};
\draw[wire] (d5.north)--(lmout.south);
\draw[wire] (sourcefinal.east)--(80,128.35)--(80,75.65)--(d2.west);
\node[nn annotation] (keys) at (87,81.6) {K、V};
\node[nn annotation,anchor=east] (query) at (123,67.15) {Q};
\node[] (pred) at (127,150.45) {cat／chases／dog};
\draw[wire] (lmout.north)--(pred.south);
\end{tikzpicture}
